\section{Results and Diagnostics}
\label{sec:results}

\subsection{Main Same-Candidate Comparison}

Table~\ref{tab:main_results} reports the main result. On all four domains,
\method{} improves over the strongest official baseline by NDCG@10: +29.7\%
on Sports, +32.2\% on Toys, +41.0\% on Home, and +43.3\% on Tools. The same
pattern holds for HR@5/@10/@20, NDCG@5/@10/@20, and MRR.

The complete evidence ledger contains 36 method-domain rows: eight official
baseline rows and one \method{} row for each of the four domains. Every row has
10,000 users, 101 candidates per user, score coverage 1.0, and the expected
row counts. For each domain, the paired test family contains 56 comparisons
(eight baselines times seven metrics); all 56 deltas are positive and
Holm-significant in every domain. Table~\ref{tab:significance_summary}
summarizes the paired-test gate, and Table~\ref{tab:full_official_ndcg10}
shows the complete per-domain NDCG@10 ranking over all eight official
baselines plus \method{}.

This supports a controlled reranking claim. It does not establish full-catalog
recommendation SOTA because candidate generation is not evaluated here.

\begin{table}[t]
\centering
\small
\caption{Paired-test gate summary for C-CRP versus the eight official baselines. The Holm family contains seven metrics for each of eight baselines per domain.}
\label{tab:significance_summary}
\begin{tabular}{lrrrr}
\toprule
Domain & Positive Holm & Min $\Delta$ & Min CI low & Max Holm $p$ \\
\midrule
Sports & 56/56 & 0.0272 & 0.0164 & 1.22e-06 \\
Toys & 56/56 & 0.0193 & 0.0083 & 4.06e-04 \\
Home & 56/56 & 0.0336 & 0.0216 & 1.02e-08 \\
Tools & 56/56 & 0.0294 & 0.0173 & 9.72e-07 \\
\bottomrule
\end{tabular}
\end{table}

\begin{table*}[t]
\centering
\scriptsize
\setlength{\tabcolsep}{3pt}
\caption{Complete same-candidate NDCG@10 ranking across all eight domains (left:
original four; right: new four). Each domain contains one \method{} row and all
eight official-code-level baseline rows over the same 101 candidates per user
(10,000 users; Beauty 973). \method{} ranks first in six of eight domains
(Books, Electronics, Sports, Toys, Home, Tools) and is competitive but not first
in the other two (Beauty, rank~2; Movies, rank~5); these two cases are shown
honestly rather than dropped. A \textit{Random} reference (analytical
expectation for one positive among 101 candidates: NDCG@10$=$0.045,
HR@10$=$0.099, MRR$=$0.051), shown for every domain, calibrates the absolute
scale; a measured \textit{Popularity} floor is additionally available for the
four new domains. Rows near this floor (e.g.\ ELMRec) retain little
same-candidate signal once their native serving stack is replaced by the unified
Qwen3-8B backbone.}
\label{tab:full_official_ndcg10}
\begin{minipage}[t]{0.49\textwidth}
\centering
\begin{tabular}{llrrrr}
\toprule
Domain & Method & Rank & NDCG@10 & HR@10 & MRR \\
\midrule
Beauty & ProEx & 1 & \textbf{0.1506} & \textbf{0.2528} & \textbf{0.1429} \\
Beauty & \method{} & 2 & 0.1341 & 0.2292 & 0.1276 \\
Beauty & IRLLRec & 3 & 0.1288 & 0.2199 & 0.1240 \\
Beauty & ProMax & 4 & 0.1270 & 0.2138 & 0.1218 \\
Beauty & RLMRec & 5 & 0.1246 & 0.1973 & 0.1249 \\
Beauty & LLMEmb & 6 & 0.1226 & 0.2384 & 0.1103 \\
Beauty & LLM-ESR & 7 & 0.1043 & 0.2066 & 0.0969 \\
Beauty & ELMRec & 8 & 0.0789 & 0.1439 & 0.0809 \\
Beauty & LLM2Rec & 9 & 0.0516 & 0.1038 & 0.0584 \\
Beauty & \textit{Random} & -- & 0.0450 & 0.0990 & 0.0515 \\
\midrule
Books & \method{} & 1 & \textbf{0.3328} & \textbf{0.4758} & \textbf{0.3063} \\
Books & LLMEmb & 2 & 0.2737 & 0.4722 & 0.2332 \\
Books & LLM2Rec & 3 & 0.2234 & 0.3720 & 0.1984 \\
Books & IRLLRec & 4 & 0.1665 & 0.3055 & 0.1550 \\
Books & RLMRec & 5 & 0.1302 & 0.2319 & 0.1270 \\
Books & ProEx & 6 & 0.1301 & 0.2568 & 0.1145 \\
Books & ProMax & 7 & 0.0980 & 0.1713 & 0.0971 \\
Books & LLM-ESR & 8 & 0.0723 & 0.1479 & 0.0732 \\
Books & ELMRec & 9 & 0.0523 & 0.1092 & 0.0578 \\
Books & \textit{Random} & -- & 0.0450 & 0.0990 & 0.0515 \\
\midrule
Elec. & \method{} & 1 & \textbf{0.1833} & \textbf{0.2993} & \textbf{0.1683} \\
Elec. & LLMEmb & 2 & 0.1196 & 0.2450 & 0.1067 \\
Elec. & IRLLRec & 3 & 0.0997 & 0.1883 & 0.0975 \\
Elec. & RLMRec & 4 & 0.0818 & 0.1612 & 0.0813 \\
Elec. & LLM-ESR & 5 & 0.0678 & 0.1426 & 0.0683 \\
Elec. & ProEx & 6 & 0.0631 & 0.1305 & 0.0665 \\
Elec. & ProMax & 7 & 0.0619 & 0.1232 & 0.0657 \\
Elec. & LLM2Rec & 8 & 0.0532 & 0.1172 & 0.0581 \\
Elec. & ELMRec & 9 & 0.0489 & 0.1050 & 0.0547 \\
Elec. & \textit{Random} & -- & 0.0450 & 0.0990 & 0.0515 \\
\midrule
Movies & LLMEmb & 1 & \textbf{0.1690} & \textbf{0.3336} & \textbf{0.1443} \\
Movies & LLM2Rec & 2 & 0.1464 & 0.2669 & 0.1324 \\
Movies & IRLLRec & 3 & 0.1397 & 0.2630 & 0.1295 \\
Movies & RLMRec & 4 & 0.1283 & 0.2471 & 0.1158 \\
Movies & \method{} & 5 & 0.1281 & 0.2083 & 0.1271 \\
Movies & ProEx & 6 & 0.1254 & 0.2450 & 0.1123 \\
Movies & ProMax & 7 & 0.1179 & 0.2184 & 0.1091 \\
Movies & LLM-ESR & 8 & 0.0745 & 0.1617 & 0.0731 \\
Movies & ELMRec & 9 & 0.0667 & 0.1408 & 0.0678 \\
Movies & \textit{Random} & -- & 0.0450 & 0.0990 & 0.0515 \\
\bottomrule
\end{tabular}
\end{minipage}
\hfill
\begin{minipage}[t]{0.49\textwidth}
\centering
\begin{tabular}{llrrrr}
\toprule
Domain & Method & Rank & NDCG@10 & HR@10 & MRR \\
\midrule
Sports & \method{} & 1 & \textbf{0.2329} & \textbf{0.3819} & \textbf{0.2076} \\
Sports & LLMEmb & 2 & 0.1795 & 0.3384 & 0.1539 \\
Sports & IRLLRec & 3 & 0.1269 & 0.2215 & 0.1244 \\
Sports & RLMRec & 4 & 0.1000 & 0.1879 & 0.0972 \\
Sports & LLM2Rec & 5 & 0.0957 & 0.2060 & 0.0883 \\
Sports & LLM-ESR & 6 & 0.0758 & 0.1564 & 0.0751 \\
Sports & ProEx & 7 & 0.0742 & 0.1527 & 0.0743 \\
Sports & ProMax & 8 & 0.0722 & 0.1387 & 0.0741 \\
Sports & ELMRec & 9 & 0.0484 & 0.1054 & 0.0537 \\
Sports & \textit{Random} & -- & 0.0450 & 0.0990 & 0.0515 \\
Sports & \textit{Pop.} & -- & 0.0484 & 0.1050 & 0.0535 \\
\midrule
Toys & \method{} & 1 & \textbf{0.2708} & \textbf{0.3964} & \textbf{0.2503} \\
Toys & LLMEmb & 2 & 0.2049 & 0.3505 & 0.1814 \\
Toys & LLM2Rec & 3 & 0.1789 & 0.3172 & 0.1592 \\
Toys & IRLLRec & 4 & 0.1338 & 0.2293 & 0.1312 \\
Toys & RLMRec & 5 & 0.1065 & 0.1885 & 0.1058 \\
Toys & ProEx & 6 & 0.0810 & 0.1615 & 0.0812 \\
Toys & ProMax & 7 & 0.0794 & 0.1435 & 0.0818 \\
Toys & LLM-ESR & 8 & 0.0546 & 0.1172 & 0.0584 \\
Toys & ELMRec & 9 & 0.0486 & 0.1043 & 0.0543 \\
Toys & \textit{Random} & -- & 0.0450 & 0.0990 & 0.0515 \\
Toys & \textit{Pop.} & -- & 0.0465 & 0.0996 & 0.0525 \\
\midrule
Home & \method{} & 1 & \textbf{0.1324} & \textbf{0.2264} & \textbf{0.1259} \\
Home & LLMEmb & 2 & 0.0939 & 0.1856 & 0.0901 \\
Home & IRLLRec & 3 & 0.0709 & 0.1443 & 0.0742 \\
Home & RLMRec & 4 & 0.0599 & 0.1268 & 0.0640 \\
Home & LLM-ESR & 5 & 0.0554 & 0.1163 & 0.0597 \\
Home & ProEx & 6 & 0.0549 & 0.1177 & 0.0593 \\
Home & LLM2Rec & 7 & 0.0509 & 0.1101 & 0.0564 \\
Home & ProMax & 8 & 0.0469 & 0.1019 & 0.0535 \\
Home & ELMRec & 9 & 0.0460 & 0.1021 & 0.0520 \\
Home & \textit{Random} & -- & 0.0450 & 0.0990 & 0.0515 \\
Home & \textit{Pop.} & -- & 0.0454 & 0.0977 & 0.0508 \\
\midrule
Tools & \method{} & 1 & \textbf{0.1661} & \textbf{0.2696} & \textbf{0.1559} \\
Tools & LLMEmb & 2 & 0.1159 & 0.2257 & 0.1065 \\
Tools & IRLLRec & 3 & 0.0853 & 0.1651 & 0.0867 \\
Tools & LLM2Rec & 4 & 0.0815 & 0.1625 & 0.0810 \\
Tools & RLMRec & 5 & 0.0684 & 0.1354 & 0.0722 \\
Tools & LLM-ESR & 6 & 0.0606 & 0.1270 & 0.0633 \\
Tools & ProEx & 7 & 0.0557 & 0.1177 & 0.0607 \\
Tools & ProMax & 8 & 0.0503 & 0.1046 & 0.0565 \\
Tools & ELMRec & 9 & 0.0459 & 0.1010 & 0.0524 \\
Tools & \textit{Random} & -- & 0.0450 & 0.0990 & 0.0515 \\
Tools & \textit{Pop.} & -- & 0.0415 & 0.0900 & 0.0486 \\
\bottomrule
\end{tabular}
\end{minipage}
\end{table*}


\subsection{Uncertainty as a Reliability Signal}

The observation module asks whether \method{} uncertainty is informative about
ranking reliability before using it as a decision feature. We compute
event-level uncertainty from the candidate-level C-CRP rows, bin events into
five uncertainty quantiles, and compare the lowest and highest bins under the
same ranking records used by the main metric importer. The package audit checks
1,010,000 finite uncertainty rows per domain, exactly 101 candidates per event,
exact candidate-key alignment, and score coverage 1.0.

Table~\ref{tab:uncertainty_stratification} reports the high-minus-low
differences for \method{}. Across all four domains, the high-uncertainty bin is
worse on NDCG@10, MRR, and HR@10. The largest gap appears on Toys
($\Delta$NDCG@10 = -0.3038), while Home and Tools still show clear NDCG@10
drops of -0.0690 and -0.0849. This supports the motivation that the uncertainty
signal stratifies same-candidate ranking reliability.

The module is descriptive. It does not show that manipulating uncertainty
causes a metric change, and it is not a substitute for the official-baseline
comparison in Table~\ref{tab:main_results}.

\begin{table}[t]
\centering
\small
\caption{C-CRP event-level uncertainty stratifies same-candidate reliability.
Values are high-uncertainty minus low-uncertainty bin differences. Negative
values mean the high-uncertainty bin performs worse.}
\label{tab:uncertainty_stratification}
\begin{tabular}{lrrr}
\toprule
Domain & $\Delta$NDCG@10 & $\Delta$MRR & $\Delta$HR@10 \\
\midrule
Sports & -0.1862 & -0.1729 & -0.2010 \\
Toys & -0.3038 & -0.2889 & -0.3135 \\
Home & -0.0690 & -0.0615 & -0.0815 \\
Tools & -0.0849 & -0.0798 & -0.0925 \\
\bottomrule
\end{tabular}
\end{table}


\subsection{Component Ablations}

Table~\ref{tab:module_evidence} summarizes the paper-critical diagnostic
modules. The component ablation is especially important because it prevents a
too-clean story. Using removal-minus-full deltas on NDCG@10, boundary
uncertainty is exactly inert across all four domains (mean delta 0). Removing
counterevidence improves the mean by 0.00168 and is nonworse in 4/4 domains;
removing the risk penalty improves the mean by 0.00082 and is nonworse in 4/4
domains. Removing calibration gap is mixed (mean -0.00078; nonworse in 2/4
domains), and removing evidence support is directionally harmful but small
(mean -0.00013; nonworse in 1/4 domains).

The honest conclusion is that the full \method{} design is a useful
paper-facing method line under the same-candidate comparison, while several
individual uncertainty components are weak or redundant under the frozen
configuration. We therefore report component ablations as supplementary
diagnostics, not as proof that every component is necessary.

\subsection{Hyperparameter Stability}

The hyperparameter module evaluates two real controls: the risk exponent
$\eta$ and the uncertainty-weight grid. Validation-selected values are compared
against reporting-only test sweeps on NDCG@10. Both controls are stable within
the pre-registered 5\% relative-drop tolerance in all four domains. For
$\eta$, the validation-selected value matches the test-best value in Sports,
Home, and Tools; Toys selects $\eta=0.25$ on validation while $\eta=0$ is
test-best, but the relative drop is only $1.25\times 10^{-5}$. For the weight
grid, the selected value is test-best on Home and within tolerance on Sports,
Toys, and Tools; the maximum relative drop is $2.97\times 10^{-4}$.

The same result also limits the interpretation. The test-best $\eta$ is zero
in all four domains, so the paper must not claim that a positive risk penalty
is necessary or uniformly beneficial. The correct claim is stability and
sensitivity, not universal optimality.

\begin{table*}[t]
\centering
\small
\caption{Paper-critical diagnostic modules. These results constrain the
mechanistic interpretation of the main comparison and are not additional
main-table SOTA evidence.}
\label{tab:module_evidence}
\begin{tabular}{p{0.18\textwidth}p{0.33\textwidth}p{0.37\textwidth}}
\toprule
Module & Supported claim & Claim boundary \\
\midrule
Uncertainty motivation
& High-uncertainty \method{} event bins have lower NDCG@10, MRR, and HR@10 than low-uncertainty bins in all four domains.
& Descriptive stratification only; no causal or statistically significant bin-effect claim. \\
\midrule
Component ablation
& Leave-one-component-out diagnostics are complete across four domains.
Evidence support is directionally helpful but small; calibration gap is mixed.
& Boundary uncertainty is inert; removing counterevidence or risk penalty is nonworse or better on NDCG@10 in all four domains. Do not claim every component is necessary. \\
\midrule
Hyperparameter stability
& For NDCG@10, $\eta$ and uncertainty-weight choices are stable in 4/4 domains within the pre-registered relative-drop tolerance.
& Test sweeps are reporting-only. Do not claim universal optima, all-metric robustness, or that positive $\eta$ is necessary. \\
\bottomrule
\end{tabular}
\end{table*}

